%==============================================================
% ΚΕΦΑΛΑΙΟ 5 — ΠΡΟΤΑΣΗ ΠΛΑΙΣΙΟΥ ΕΠΙΚΥΡΩΣΗΣ
%==============================================================
\chapter{Πρόταση Πλαισίου Επικύρωσης}
\label{ch:framework}

Όσα προηγήθηκαν στην εργασία συνδυάζονται στο παρόν κεφάλαιο για
την δημιουργία και διατύπωση μιας σειράς προτάσεων για ένα πλαίσιο
επικύρωσης μοντέλων μηχανικής μάθησης στον πιστωτικό κίνδυνο.
Τα ευρήματα του Κεφαλαίου~\ref{ch:results} έχουν σημαντικό ρόλο
στη διαμόρφωσή του. Η ενότητα~\ref{sec:principles} αφορά τις
βασικές αρχές σχεδιασμού του πλαισίου. Στην
ενότητα~\ref{sec:workflow} περιγράφεται η πρότυπη ροή εργασιών
επικύρωσης μοντέλων μηχανικής μάθησης. Η
ενότητα~\ref{sec:complexity_requirements} έχει ως κύριο θέμα τις
διαφοροποιήσεις λόγω της πολυπλοκότητας κάθε μοντέλου και τέλος
η ενότητα~\ref{sec:limitations} περιέχει τους περιορισμούς και
τις γενικεύσεις του πλαισίου.

%==============================================================
\section{Αρχές Σχεδιασμού}
\label{sec:principles}
%==============================================================

Για τη δόμηση του πλαισίου επικύρωσης είναι αναγκαία η διαίρεση
του σε βασικές αρχές. Η ανάλυση της θεωρίας στο
Κεφάλαιο~\ref{ch:theory} και τα πειραματικά αποτελέσματα του
Κεφαλαίου~\ref{ch:results} οδηγούν στις ακόλουθες αναγκαίες
αρχές, θεμέλια για το προτεινόμενο πλαίσιο.

\paragraph{Αρχή 1 --- Αναλογικότητα ελέγχου με την πολυπλοκότητα.}
Η εντατικότητα του ελέγχου πρέπει να είναι αναλογική όπου χρειάζεται,
 με βάση εμπειρική ανάλυσης, της
πολυπλοκότητας του μοντέλου. Η παρούσα αρχή ικανοποιεί τις απαιτήσεις των EBA Guidelines και του 
SR Letter~11-7 \citep{eba2023, sr117}, οι οποίες υιοθετούν την αρχή της 
proportionality στο model risk management \citep{bis2024principles}. Η ανάλυση
του μοντέλου FFNN απαίτησε ${\sim}2{,}5$ ώρες για τον υπολογισμό
τιμών SHAP ενώ το XGBoost μερικά δευτερόλεπτα
(βλ.\ ενότητα~\ref{sec:results_ffnn_shap}).

\paragraph{Αρχή 2 --- Η ερμηνευσιμότητα ως μη διαπραγματεύσιμη απαίτηση.}
Για κάθε μοντέλο ανεξαρτήτως πολυπλοκότητας είναι αναγκαία η
ύπαρξη μηχανισμού επεξήγησης των αποτελεσμάτων του σε ατομικό
αλλά και καθολικό επίπεδο. Οι ρυθμιστικές αρχές υποχρεώνουν
την ύπαρξη ανάλογου μηχανισμού για την ικανοποίηση του
«δικαιώματος εξήγησης» \citep{goodman2017european, eba2023} σε πρώτο ή και δεύτερο επίπεδο, όπως
αναφέρθηκε στην ενότητα~\ref{sec:results_pdp}. Η επιλογή
μοντέλου λοιπόν πρέπει να περιορίζεται σε αλγορίθμους οι οποίοι
συνοδεύονται από αξιόπιστο μηχανισμό post-hoc εξήγησης. Από τα
αποτελέσματα της ενότητας~\ref{sec:results_ffnn_shap} φαίνεται
πρακτικά ότι η ύπαρξη ενός τέτοιου μηχανισμού και η αξιοπιστία
του εξαρτάται από την πολυπλοκότητα του μοντέλου, συνδέοντας
τις Αρχές 1 και 2.

\paragraph{Αρχή 3 --- Διαχωρισμός μεροληψίας μοντέλου από μεροληψία δεδομένων.}
Πρέπει να γίνεται έρευνα για ανίχνευση μεροληψίας σε όλα τα
μοντέλα ώστε να γίνεται εμφανής η πηγή της. Μεροληψία μπορεί
να εμφανιστεί λόγω της βάσης δεδομένων
(βλ.\ ενότητα~\ref{subsec:bias_detection}) αλλά να ερμηνευθεί
λανθασμένα ως σφάλμα στην εσωτερική λογική κάποιου μοντέλου.
Για την εμπειρική απόδειξη της πηγής της μεροληψίας είναι
απαραίτητη η σύγκριση των αποτελεσμάτων των μοντέλων. Η
συγκεκριμένη αρχή προσθέτει στη διαδικασία μετριασμού
μεροληψίας, εφόσον δεν στοχεύει μόνο στον μετριασμό της αλλά
και στην αιτία εμφάνισής της.

\paragraph{Αρχή 4 --- Παρακολούθηση ως συνεχής διαδικασία και όχι ως τελικός έλεγχος.}
Η διαδικασία της επικύρωσης πρέπει να είναι μια διαρκής
διαδικασία. Η μεταβολή στα δεδομένα (Data Drift) και στο
μοντέλο (Model Drift) έχει αποδειχθεί από τη βιβλιογραφία ότι
είναι αναπόφευκτη~\citep{widmer1996learning}. Οι αλγόριθμοι
θα πρέπει να παρακολουθούνται περιοδικά με δείκτες όπως ο PSI
και CSI, οι οποίοι όπως απέδειξε η εργασία είναι ικανοί να
παράγουν αξιόπιστες προειδοποιήσεις. Επιτρέπεται έτσι η
διάγνωση και η έγκαιρη επανεκπαίδευση του μοντέλου προς την
αποφυγή αισθητής υποβάθμισης των προβλέψεων.


\paragraph{Αρχή 5 --- Κανονιστική συμμόρφωση ως ελάχιστο, όχι ως στόχος.}
Οι απαραίτητες προς συμμόρφωση απαιτήσεις των Basel~II/III,
SR~11-7 και EBA Guidelines \citep{basel2017, sr117, eba2023} αποτελούν ελάχιστο επίπεδο προς
κάλυψη και όχι όριο. Για τους δείκτες DPD, DI, καθώς και PSI,
CSI, νομοθετούνται όρια των οποίων η ικανοποίηση δεν εγγυάται
την πλήρη ηθική κάλυψη των αποφάσεων του μοντέλου. Είναι
γνωστό από τη βιβλιογραφία ότι η νομική συμμόρφωση δεν
εξαντλεί την ηθική υποχρέωση~\citep{barocas2019fairness, selbst2019fairness},
ακόμη και όταν η ηθική κάλυψη συνεπάγεται απώλεια
προβλεπτικής ικανότητας. Η αρνητική αυτή επίπτωση κοστίζει
στον τραπεζικό τομέα, όμως η απώλεια μπορεί να μετριαστεί
και να διατηρηθεί σε χαμηλά επίπεδα. Στην εργασία φτάνουμε
σε τέλεια βαθμολογία DI $= 1$ με Balanced Accuracy μόλις
$-0{,}022$.

%==============================================================
\section{Βήματα Validation Workflow}
\label{sec:workflow}
%==============================================================

Για τη ροή επικύρωσης προτείνεται διαδικασία επτά βημάτων.
Τα βήματα πρέπει να εφαρμόζονται με τη σειρά που δίνονται.
Συγκεκριμένα βήματα διαφοροποιούνται ανάλογα με την
πολυπλοκότητα του μοντέλου, όπως φαίνεται αναλυτικά στην
ενότητα~\ref{sec:complexity_requirements}.

\paragraph{Βήμα 1 --- Προεπεξεργασία και διαχωρισμός δεδομένων.}
Εάν η βάση περιέχει εγγραφές χωρίς πληροφορία, τότε πρέπει
οι εγγραφές αυτές είτε να διαγράφονται ή να συμπληρώνονται
με εκτιμώμενες τιμές ανάλογα με τη βάση. Κωδικοποίηση
κατηγορικών μεταβλητών ώστε να μπορούν να χρησιμοποιηθούν
από αλγορίθμους μηχανικής μάθησης. Στη συνέχεια ο
διαχωρισμός σε train, validation και test sets πρέπει να
γίνεται με stratification ώστε να εξασφαλίζεται η αναλογία
στις κλάσεις. Τέλος, εάν απαιτείται κανονικοποίηση, αυτή
πρέπει να υλοποιείται με fit μόνο στο training set, ενώ στα
test και validation sets με transform, για να αποφευχθεί
το data leakage.

\paragraph{Βήμα 2 --- Εκπαίδευση μοντέλου και έλεγχος υπερπροσαρμογής.}
Μετά την επιλογή υπερπαραμέτρων (βάθος δέντρου, learning
rate, regularization strength) και την εκπαίδευση με αυτές,
πρέπει να γίνεται άμεσα σύγκριση απόδοσης στα train και
validation sets και έλεγχος υπερπροσαρμογής. Ως ένδειξη
υπερπροσαρμογής μπορεί να χρησιμοποιηθεί η σχετική πτώση
μεγαλύτερη του ${\sim}5\%$ (ad-hoc όριο βασισμένο στα ευρήματα
του Κεφαλαίου~\ref{ch:results}). Σε περίπτωση υπερπροσαρμογής
πρέπει να γίνει επανεξέταση των παραμέτρων ή χρήση τεχνικών
regularization (βλ.\ Πίνακα~\ref{tab:overfitting}). Η
οριακή παραβίαση από μοντέλα δέντρων αποφάσεων μπορεί να
αντιμετωπιστεί με early stopping και stochastic
regularization, οπότε επιτρέπεται μικρή ανοχή για την
περίπτωση αυτή.

\paragraph{Βήμα 3 --- Αξιολόγηση απόδοσης σε πολλαπλούς άξονες.}
Υπολογισμός των μετρικών AUC-ROC, Gini και KS στο test set,
καθώς και ανάλυση πίνακα σύγχυσης, καμπυλών Precision/Recall
ως προς το threshold, ανάλυση κατά δεκατημόρια και
διαγραμμάτων βαθμονόμησης. Κάθε μετρική είναι απαραίτητη
και συμπληρώνει την εικόνα για τη λειτουργία του μοντέλου.
Επιπλέον είναι αναγκαίες για τη συμμόρφωση με τις εποπτικές
αρχές~\citep{basel2017, sr117, eba2023}.

\paragraph{Βήμα 4 --- Ερμηνεία αποφάσεων.}
Ερμηνεία αποφάσεων μοντέλου μέσω SHAP feature importance και
beeswarm διαγραμμάτων για καθολική σημαντικότητα μεταβλητών,
καθώς και SHAP waterfall plots και Partial Dependence Plots
για ατομική σημαντικότητα χαρακτηριστικών ανά
μοντέλο~\citep{lundberg2017shap}. Για μοντέλα δέντρων
αποφάσεων είναι προτιμότερη η συνάρτηση TreeExplainer
(ακριβής, υπολογιστικά αποδοτική), ενώ για νευρωνικά δίκτυα
η KernelExplainer, η οποία όμως επιφέρει σημαντικό
υπολογιστικό κόστος.

\paragraph{Βήμα 5 --- Ανίχνευση μεροληψίας σε όλα τα μοντέλα.}
Αρχικά πρέπει να γίνει επιλογή ευαίσθητου χαρακτηριστικού
και διαχωρισμός των δεδομένων σε ομάδες βάσει αυτού.
Υπολογισμός DPD και DI με τη βιβλιοθήκη
Fairlearn~\citep{bird2020fairlearn} σε πολλαπλά μοντέλα,
ώστε σε πιθανή σύγκλιση αποτελεσμάτων να εντοπιστεί ότι η
βάση αποτελεί την πηγή μεροληψίας.

\paragraph{Βήμα 6 --- Μετριασμός μεροληψίας με αξιολόγηση κόστους.}
Εφόσον εντοπιστεί μεροληψία, εφαρμόζεται τεχνική μετριασμού
(in-processing για Logistic Regression, post-processing για
XGBoost). Η τεχνική μετριασμού θα επιφέρει κόστος στην
απόδοση, το οποίο πρέπει να καταγραφεί. Στην εργασία κατά
το post-processing (ThresholdOptimizer) διατηρείται το AUC
σταθερό αλλά μειώνεται το Balanced Accuracy, ενώ κατά το
in-processing (\texttt{ExponentiatedGradient}) μειώνεται η
AUC κατά ${\sim}14\%$. Η εξισορρόπηση ανάμεσα στην
αξιοπιστία πρόβλεψης έναντι της ισορροπίας ομάδων
εξαρτάται από τους ισχύοντες νόμους, την ηθική πυξίδα και
τις επιχειρηματικές προτεραιότητες της τράπεζας.

\paragraph{Βήμα 7 --- Συνεχής παρακολούθηση μετά το deployment.}
Συνεχής υπολογισμός PSI και CSI για πιθανή αλλαγή στα
δεδομένα εισόδου και τις προβλέψεις των
αλγορίθμων~\citep{eba2023}. Συγκεκριμένα, τιμές κάτω του
$0{,}10$ δηλώνουν σταθερότητα, μεταξύ $0{,}10$ και $0{,}25$
απαιτούν αυξημένη παρακολούθηση και άνω του $0{,}25$
επιβάλλουν επανεκπαίδευση. Σε περιπτώσεις αστάθειας
απαιτείται ανάλυση SHAP για τον εντοπισμό των
χαρακτηριστικών που την προκάλεσαν. Όπως αναφέρθηκε στο
Βήμα~4, το υπολογιστικό κόστος της ανάλυσης διαφέρει
σημαντικά ανά μοντέλο και πρέπει να ληφθεί υπόψη εκ των
προτέρων.

%==============================================================
\section{Απαιτήσεις Επικύρωσης ανά Επίπεδο Πολυπλοκότητας}
\label{sec:complexity_requirements}
%==============================================================

Οι παρακάτω Πίνακες~\ref{tab:complexity_a} και~\ref{tab:complexity_b}
παρουσιάζουν τις διαφορές στις απαιτήσεις επικύρωσης για τα τρία
μοντέλα της εργασίας. Μέσω των αποτελεσμάτων που παρουσιάζονται
γίνεται σαφής η διαφοροποίηση στο πλαίσιο επικύρωσης ανάλογα με
την πολυπλοκότητα του αλγορίθμου.

\begin{table}[htbp]
    \centering
    \renewcommand{\arraystretch}{1.3}
    \caption[Απαιτήσεις επικύρωσης ανά πολυπλοκότητα — Ερμηνευσιμότητα \& Απόδοση]{Διαφοροποιημένες απαιτήσεις επικύρωσης ανά επίπεδο
             πολυπλοκότητας --- Ερμηνευσιμότητα \& Απόδοση.
             Οι τιμές προέρχονται από τα ευρήματα του
             Κεφαλαίου~\ref{ch:results}.}
    \label{tab:complexity_a}
    \small
    \begin{tabular}{p{3.8cm}p{3.8cm}p{3.8cm}}
        \toprule
        \textbf{Απαίτηση} &
        \textbf{Logistic Regression} &
        \textbf{XGBoost} \\
        \midrule
        Εγγενής ερμηνευσιμότητα &
        Πλήρης (συντελεστές $\beta$) &
        Καμία \\
        \midrule
        Post-hoc εξήγηση &
        Δεν απαιτείται &
        SHAP TreeExplainer (δευτερόλεπτα) \\
        \midrule
        Έλεγχος υπερπροσαρμογής &
        Σχεδόν περιττός (1{,}3\%) &
        Κρίσιμος (6{,}2\%) \\
        \midrule
        Αξιοπιστία βαθμονόμησης &
        Υψηλή (log-loss training) &
        Μέτρια (απόκλιση στις άκρες) \\
        \bottomrule
    \end{tabular}

    \vspace{0.4cm}

    \begin{tabular}{p{3.8cm}p{3.8cm}p{3.8cm}}
        \toprule
        \textbf{Απαίτηση} &
        \textbf{XGBoost} &
        \textbf{FFNN} \\
        \midrule
        Εγγενής ερμηνευσιμότητα &
        Καμία &
        Καμία \\
        \midrule
        Post-hoc εξήγηση &
        SHAP TreeExplainer (δευτερόλεπτα) &
        SHAP KernelExplainer (${\sim}2{,}5$ ώρες) \\
        \midrule
        Έλεγχος υπερπροσαρμογής &
        Κρίσιμος (6{,}2\%) &
        Μέτριος (2{,}6\%) \\
        \midrule
        Αξιοπιστία βαθμονόμησης &
        Μέτρια (απόκλιση στις άκρες) &
        Χαμηλή (ακανόνιστη) \\
        \bottomrule
    \end{tabular}
\end{table}

\begin{table}[htbp]
    \centering
    \renewcommand{\arraystretch}{1.3}
    \caption[Απαιτήσεις επικύρωσης ανά πολυπλοκότητα — Δικαιοσύνη \& Παρακολούθηση]{Διαφοροποιημένες απαιτήσεις επικύρωσης ανά επίπεδο
             πολυπλοκότητας --- Δικαιοσύνη \& Παρακολούθηση.
             Οι τιμές προέρχονται από τα ευρήματα του
             Κεφαλαίου~\ref{ch:results}.}
    \label{tab:complexity_b}
    \small
    \begin{tabular}{p{3.8cm}p{3.8cm}p{3.8cm}}
        \toprule
        \textbf{Απαίτηση} &
        \textbf{Logistic Regression} &
        \textbf{XGBoost} \\
        \midrule
        Τεχνική μετριασμού bias &
        In-processing (\texttt{Exp.\ Gradient}) &
        Post-processing (Thresh.\ Opt.) \\
        \midrule
        Κόστος μετριασμού bias &
        AUC: $-14\%$ &
        AUC: $0\%$ / Bal.\ Acc.: $\downarrow$ \\
        \midrule
        Παρακολούθηση (PSI/CSI) &
        Τυπική &
        Τυπική \\
        \midrule
        Σύνδεση drift με χαρακτηριστικά &
        Άμεση (συντελεστές) &
        Γρήγορη (TreeExplainer) \\
        \midrule
        Συνολικό κόστος επικύρωσης &
        Χαμηλό &
        Μέτριο \\
        \bottomrule
    \end{tabular}

    \vspace{0.4cm}

    \begin{tabular}{p{3.8cm}p{3.8cm}p{3.8cm}}
        \toprule
        \textbf{Απαίτηση} &
        \textbf{XGBoost} &
        \textbf{FFNN} \\
        \midrule
        Τεχνική μετριασμού bias &
        Post-processing (Thresh.\ Opt.) &
        Δεν εξετάστηκε \\
        \midrule
        Κόστος μετριασμού bias &
        AUC: $0\%$ / Bal.\ Acc.: $\downarrow$ &
        --- \\
        \midrule
        Παρακολούθηση (PSI/CSI) &
        Τυπική &
        Τυπική (CSI ταυτόσημο) \\
        \midrule
        Σύνδεση drift με χαρακτηριστικά &
        Γρήγορη (TreeExplainer) &
        Πρακτικά απαγορευτική (KernelExplainer) \\
        \midrule
        Συνολικό κόστος επικύρωσης &
        Μέτριο &
        Πολύ υψηλό \\
        \bottomrule
    \end{tabular}
\end{table}


Τα κύρια συμπεράσματα είναι τρία. Αρχικά, οι απαιτήσεις επικύρωσης
δεν αυξάνονται γραμμικά όσο αυξάνεται η πολυπλοκότητα αλλά
δυσανάλογα περισσότερο. Πιο συγκεκριμένα, η διαφορά ανάμεσα στις
απαιτήσεις για το μοντέλο XGBoost και την Logistic Regression είναι
σημαντικά μικρότερη από τη διαφορά XGBoost και FFNN. Ταυτόχρονα η
τελευταία αύξηση στην πολυπλοκότητα από XGBoost σε FFNN δεν εισάγει
ουσιαστικό πλεονέκτημα στην απόδοση του αλγορίθμου. Δεύτερον, οι
απαιτήσεις παρακολούθησης του μοντέλου παραμένουν σταθερές, αλλά η
ανάγκη ερμηνείας τους εμφανίζει διαφοροποίηση στο κόστος υπολογισμού.
Η παρακολούθηση μοντέλου χρησιμοποιεί τις μετρικές PSI και CSI που
υπολογίζουν διαφορές στα δεδομένα που εισάγονται στα μοντέλα, άρα ο
υπολογισμός τους δεν διαφέρει ανά μοντέλο. Σε περίπτωση όμως
διαπίστωσης σημαντικής αλλαγής στα δεδομένα εισαγωγής είναι αναγκαία
η ερμηνεία τους μέσω SHAP που διαφέρει ανά μοντέλο. Συμπεραίνεται
έτσι ότι η απόδοση δεν είναι το μοναδικό κριτήριο για την επιλογή
του μοντέλου, καθώς η επικύρωση και η συνεχής παρακολούθησή του,
οι οποίες είναι αναγκαίες, αποτελούν παράγοντες βιωσιμότητάς του.

Με βάση τα παραπάνω, για τη συγκεκριμένη βάση το μοντέλο XGBoost
κρίνεται ως το καταλληλότερο από τα τρία μοντέλα. Η προβλεπτική
ικανότητά του υπερέχει έστω και ελαφρά έναντι των άλλων μοντέλων,
ενώ το κόστος ερμηνείας του παραμένει σε αποδεκτά όρια. Επιπλέον,
επιτρέπει τον μετριασμό μεροληψίας post-processing, ο οποίος δεν
εμφανίζει απώλειες στον δείκτη AUC, με μόλις $-0{,}022$ στο
Balanced Accuracy, και δεν απαιτεί επανεκπαίδευση μοντέλου. Το
μοντέλο FFNN εμφανίζει καλύτερη απόδοση έναντι του XGBoost μόνο
στον δείκτη KS ($0{,}465$ έναντι $0{,}458$) με αποτέλεσμα να μην
δικαιολογείται το δυσανάλογο κόστος ερμηνείας του.

%==============================================================
\section{Περιορισμοί και Γενικευσιμότητα}
\label{sec:limitations}
%==============================================================

Το πλαίσιο που προτείνεται, παρότι προκύπτει από τα αποτελέσματα
και το κανονιστικό πλαίσιο, χαρακτηρίζεται από περιορισμούς.
Η ανάλυση αυτών των περιορισμών κρίνεται αναγκαία, καθώς
επηρεάζεται άμεσα η γενικευσιμότητα των πορισμάτων.

\paragraph{Ένα μόνο dataset.}
Στην παρούσα εργασία χρησιμοποιείται μόνο μια βάση δεδομένων,
η FICO HELOC. Η συγκεκριμένη βάση έχει μερικές ιδιαιτερότητες.
Αρχικά, περιέχει συγκεντρωτικά χαρακτηριστικά όπως το
\texttt{ExternalRiskEstimate}. Τα συγκεκριμένα χαρακτηριστικά
εμφανίζουν σχετικά μονότονη συμπεριφορά σε σχέση με τον
πιστωτικό κίνδυνο, με αποτέλεσμα η διαφορά στην απόδοση μεταξύ
της Logistic Regression --- η οποία δεν αναγνωρίζει μη γραμμικές
σχέσεις ανάμεσα στα χαρακτηριστικά και την τελική πρόβλεψη ---
και πιο σύνθετων μοντέλων να μειώνεται. Σε βάσεις δεδομένων με
μη γραμμικά χαρακτηριστικά η διαφορά απόδοσης ενδέχεται να είναι
σημαντικότερη. Επιπλέον, η βάση δεν περιέχει δημογραφικά στοιχεία.
Η απουσία τους οδήγησε στον ορισμό proxy μεταβλητής, και
συγκεκριμένα της \texttt{AverageMInFile}. Η μεταβλητή
\texttt{AverageMInFile} δεν αποτελεί προστατευόμενο χαρακτηριστικό
αλλά είναι feature κινδύνου. Λόγω αυτού, η μείωση στην απόδοση
λόγω της διαδικασίας εξάλειψης της μεροληψίας ενδέχεται να είναι
τεχνικά αυξημένη. Η επιλογή proxy μεταβλητής είναι θεμιτή ως προς
την επίδειξη της μεθοδολογίας, αλλά δεν αντικαθιστά πλήρως την
ανάλυση σε πραγματικές προστατευόμενες μεταβλητές. Συνεπώς, σε
βάσεις δεδομένων όπου η εξάλειψη μεροληψίας γίνεται σε
προστατευόμενα χαρακτηριστικά, αναμένεται η απώλεια απόδοσης
λόγω mitigation να είναι μικρότερη.

\paragraph{Στατική προσομοίωση μετατόπισης δεδομένων.}
Η διαδικασία παρακολούθησης των μοντέλων βασίστηκε σε τρία
σενάρια οικονομικής επιδείνωσης (ήπια, μέτρια, σοβαρή) τα οποία
παράχθηκαν τεχνικά με τη χρήση Γκαουσιανού θορύβου. Η προσέγγιση
αυτή, ενώ είναι κατάλληλη για επίδειξη χρήσης των δεικτών PSI και
CSI, δεν περιέχει ακριβή στοιχεία οικονομικής κρίσης. Κατά την
περίοδο μιας οικονομικής κρίσης αναμένεται να εμφανίσουν
συσχετισμένες μετατοπίσεις πολλαπλά χαρακτηριστικά, σε αντίθεση
με την προσομοίωση που εμφάνισαν μόνο τρία, ενώ στα υπόλοιπα
είχε προστεθεί απλά λευκός θόρυβος.

\paragraph{Εστίαση στην πιθανότητα αθέτησης.}
Στον τραπεζικό τομέα, για τη συνολική διαχείριση ενός
χαρτοφυλακίου είναι αναγκαία η πρόβλεψη της Ζημίας σε Περίπτωση
Αθέτησης (LGD) και της Έκθεσης σε Περίπτωση Αθέτησης (EAD),
εκτός από την Πιθανότητα Αθέτησης (PD)~\citep{mcneil2015}. Το
πλαίσιο που προτείνεται στην εργασία δεν επεκτείνεται στην
πρόβλεψη των δύο πρώτων μεταβλητών, καθώς υπερβαίνει τον σκοπό
της. Ως αποτέλεσμα, το παρόν πλαίσιο δεν αποτελεί ολοκληρωμένη
πρόταση για τη διαχείριση χαρτοφυλακίου, αλλά μόνο για την
πρόβλεψη της πιθανότητας αθέτησης.

\paragraph{Επιλογή υπερπαραμέτρων.}
Δεν έγινε βελτιστοποίηση υπερπαραμέτρων πριν από την εκπαίδευση
των μοντέλων για την επίτευξη βέλτιστης απόδοσης. Στην παρούσα
εργασία τα τρία μοντέλα εκπαιδεύτηκαν με λογικές αρχικές τιμές,
οι οποίες όμως δεν παράγουν τη βέλτιστη δυνατή απόδοση των
μοντέλων. Η επιλογή αυτή έγινε διότι, μέσω της βελτιστοποίησης
υπερπαραμέτρων, ενώ ενδέχεται να υπάρξουν ελαφρές βελτιώσεις
στην απόδοση των μοντέλων, δεν επηρεάζονται οι απαιτήσεις
ερμηνείας και παρακολούθησης. Έτσι, η εξαντλητική χρήση της
αποτελεί επιλογή για εργασίες με κεντρικό ζήτημα τη
μεγιστοποίηση της απόδοσης κάθε μοντέλου.

\paragraph{Γενικευσιμότητα του πλαισίου.}
Παρά τους περιορισμούς που αναφέρθηκαν, η δομή του πλαισίου που
προτείνεται παραμένει εφαρμόσιμη. Οι αρχές σχεδιασμού, η πρότυπη
ροή εργασιών και οι κλιμακούμενες απαιτήσεις ανά πολυπλοκότητα
επεκτείνονται σε προβλήματα πρόβλεψης PD σε όλα τα είδη δανείων,
σε υλοποιήσεις άλλων συγγενικών αλγορίθμων (LightGBM, CatBoost),
αλλά και σε άλλα binary classification προβλήματα του
χρηματοπιστωτικού τομέα (ανίχνευση απάτης, prepayment risk). Η
βασική ιδέα επικύρωσης και ερμηνευσιμότητας αναλογικά με την
πολυπλοκότητα παραμένει αναλλοίωτη, ανεξαρτήτως βάσης δεδομένων
και αλγορίθμων. Έτσι, αποτελεί σαφή απάντηση στο ερώτημα της
εργασίας.